% !TEX root = ./main.tex
\chapter*{Glossario}
\addcontentsline{toc}{chapter}{Glossario}
\label{app:glossary}

\begingroup
\small
\setlength{\parindent}{0pt}
\setlength{\parskip}{0.65em}

Questo glossario raccoglie solo i termini tecnici che tornano più spesso nella tesi e che, fuori contesto, non sono così immediati. Le voci sono ordinate alfabeticamente.

\textbf{Affordance.} Insieme di segnali visivi o interattivi che fanno intuire come un elemento dell'interfaccia possa essere usato.

\textbf{API (\emph{Application Programming Interface}).} Interfaccia con cui un componente software espone funzioni o dati ad altri componenti.

\textbf{ASGI (\emph{Asynchronous Server Gateway Interface}).} Standard usato in Python per applicazioni web asincrone; nel progetto fa da base ai servizi FastAPI.

\textbf{Backend.} Parte del sistema che gestisce logica applicativa, servizi, orchestrazione e accesso ai dati, senza coincidere con l'interfaccia mostrata all'utente.

\textbf{Blendshape.} Tecnica di animazione facciale che modifica la mesh 3D per ottenere espressioni, movimenti labiali o transizioni del volto.

\textbf{BM25 (\emph{Best Matching 25}).} Algoritmo di ranking lessicale usato per stimare quanto un testo sia pertinente rispetto a una query.

\textbf{Chunk.} Porzione di testo ricavata dalla suddivisione di un contenuto più grande, così da renderlo indicizzabile e recuperabile.

\textbf{CLI (\emph{Command-Line Interface}).} Interfaccia testuale usata per avvio, diagnostica o manutenzione di servizi e strumenti.

\textbf{CORS (\emph{Cross-Origin Resource Sharing}).} Meccanismo del browser che regola quali risorse possono essere richieste tra origini diverse.

\textbf{Debriefing.} Breve momento finale della prova in cui il partecipante può commentare liberamente l'esperienza e chiarire eventuali dubbi sulla sessione svolta.

\textbf{Embedding.} Rappresentazione numerica di un testo, di un'immagine o di un altro contenuto, usata per confrontarne la vicinanza semantica.

\textbf{ECA (\emph{Embodied Conversational Agent}).} Agente conversazionale dotato di una rappresentazione corporea visibile, come un avatar o un personaggio in scena.

\textbf{Fallback.} Comportamento di riserva adottato quando il percorso principale non può essere completato o non restituisce un risultato affidabile.

\textbf{FSM (\emph{Finite State Machine}).} Modello a stati finiti usato per descrivere passaggi ordinati tra fasi diverse di un sistema o di un'interfaccia.

\textbf{Grounding.} Ancoraggio della risposta a fonti, memoria o input reali, in modo da ridurre invenzioni e mantenere coerenza col contesto.

\textbf{Hallucination.} Generazione di un contenuto plausibile ma non supportato dai dati disponibili; nel testo il termine viene spesso reso anche come ``allucinazione''.

\textbf{HMD (\emph{Head-Mounted Display}).} Visore indossabile usato per la fruizione di esperienze VR o XR.

\textbf{Hybrid search.} Strategia di recupero che combina confronto lessicale e ricerca semantica per migliorare la pertinenza dei contenuti trovati.

\textbf{IQR (\emph{Interquartile Range}).} Intervallo tra primo e terzo quartile, usato per descrivere la dispersione dei dati in modo poco sensibile ai valori estremi.

\textbf{LLM (\emph{Large Language Model}).} Modello linguistico di grandi dimensioni impiegato per comprendere il contesto e generare risposte in linguaggio naturale.

\textbf{MDMT (\emph{Multidimensional Measure of Trust}).} Questionario usato per descrivere la fiducia percepita verso un agente o un sistema.

\textbf{OCR (\emph{Optical Character Recognition}).} Tecnica usata per estrarre testo da immagini o PDF scannerizzati, così da renderli interrogabili.

\textbf{PANAS (\emph{Positive and Negative Affect Schedule}).} Questionario che misura in modo separato affettività positiva e affettività negativa.

\textbf{Prompt.} Istruzione o insieme di istruzioni date a un modello generativo per orientarne comportamento, tono o vincoli di risposta.

\textbf{RAG (\emph{Retrieval-Augmented Generation}).} Approccio che combina recupero di contenuti pertinenti e generazione testuale, così da produrre risposte più ancorate al contesto.

\textbf{Ranking.} Ordinamento dei risultati recuperati in base alla loro pertinenza rispetto alla richiesta dell'utente.

\textbf{Retrieval.} Fase di recupero di documenti, chunk o memorie rilevanti da usare nella generazione della risposta.

\textbf{Routing.} Scelta del percorso operativo più adatto in base al tipo di richiesta o di intento rilevato nel turno corrente.

\textbf{RSES (\emph{Rosenberg Self-Esteem Scale}).} Scala usata per ottenere una misura sintetica dell'autostima percepita.

\textbf{STT (\emph{Speech-To-Text}).} Componente che trasforma il parlato in testo e apre la parte linguistica della pipeline conversazionale.

\textbf{TTS (\emph{Text-To-Speech}).} Componente che sintetizza una voce a partire da un testo generato dal sistema.

\textbf{UI (\emph{User Interface}).} Parte visibile e interattiva del sistema con cui l'utente entra in contatto diretto.

\textbf{UX (\emph{User Experience}).} Qualità complessiva dell'esperienza d'uso, intesa come chiarezza, fluidità e percezione del comportamento del sistema.

\textbf{Vector store.} Archivio specializzato che conserva embedding e metadati, rendendo possibile il retrieval semantico tramite similarità vettoriale.

\textbf{Voice cloning.} Tecnica che permette di sintetizzare una voce con caratteristiche simili a quelle di un campione vocale di riferimento.

\textbf{WebAssembly.} Formato binario eseguibile nel browser, usato per componenti web ad alte prestazioni.

\textbf{WebGL (\emph{Web Graphics Library}).} API grafica del browser che permette di eseguire contenuti 2D e 3D accelerati lato client.

\textbf{Zero-shot.} Modalità in cui un modello affronta un compito senza esempi specifici addestrati per quel caso particolare.

\endgroup
